\scanpage{99}
\mspage{14}
% archive sheet number: 98

de pureté cohomologique est vrai pour le couple $(X',Y')$ et pour $\ell$.

Sous ces conditions, \inserted{soit $U$ lisse quasi-projectif sur $K$, et} pour \unsure{tt|$\forall$} entier $i$, l'action du groupe d'inertie $I$ sur $\deleted{\unsure{H^i(U_{\bar K})}}H^i(Y_{\bar K},\mathbb Z_\ell)$ est \deleted{\unsure{quasi}}\inserted{\unsure{essentiellement}} unipotente d'échelon $i+1$.

\uline{Dém} En vertu de I\ldots, on peut supposer $S=\hat S$. \deleted{On peut que} Comme $U$ est quasi-projectif sur $K$, \inserted{de dim relative $n$,} \deleted{il provient d'un} $\exists X$ \deleted{projectif} propre sur $S$ (et = projectif sur $S$), \deleted{plat sur $S$} de dim relative $n$. On peut supposer $X=$ adhérence schématique de $U$, donc $X$ plat sur $S$ et $X$ réduit. \deleted{Notons} \inserted{Rappelons} que l'hypothèse de \unsure{lissité} \deleted{de} sur $U$ implique $H^i(U_{\bar K},\mathbb Z/\ell\mathbb Z)$ fini, donc $E=H^i(U_{\bar K},\mathbb Z_\ell)$ de t.f. sur $\mathbb Z_\ell$.

Quitte : faire une extension finie \inserted{séparable} sur $K$, on peut supposer que l'action de $P$ sur $E$ est triviale (cf I\ldots). Utilisant la résolution \deleted{des} forte des singularités \inserted{(hyp. (i))}, on est ramené au cas $X$ régulier, et $Y=X-U$ un diviseur à croisements normaux. On peut alors appliquer 3.2 grâce à l'hypothèse (ii), et on gagne.

\uline{Remarque} 3.6. Si on admet le résultat de compactification de Nagata, la démonstration précédente marche
